% Copyright (c) 2026 Eugene Joseph M. Ragasa.
% SPDX-License-Identifier: Apache-2.0
% Author: Eugene Joseph M. Ragasa
% Formal proof completed: 2026-10-19
\documentclass[11pt]{article}
\usepackage[T1]{fontenc}
\usepackage{amsmath,amssymb,booktabs}
\usepackage[margin=1in]{geometry}
\usepackage[hidelinks]{hyperref}
\usepackage{tabularx}
\title{Proof Dossier: Spectral Rectangle-Containment Transfer}
\author{Eugene Joseph M. Ragasa}
\date{Formal proof completed: 2026-10-19}
\begin{document}
\maketitle
\section{Identity and status}
\begin{tabularx}{\textwidth}{@{}l>{\raggedright\arraybackslash}X@{}}
\toprule
Statement ID & \texttt{BG-THM-031} \\
Status & \texttt{proved} \\
Lean declaration & \path{ProjectKoios.BandGauge.Geometry.rankTwoSpectralEllipseInRectangle_eq_full_of_subset} \\
Module & \path{ProjectKoios/BandGauge/Geometry/RankTwoLoss.lean} \\
\bottomrule
\end{tabularx}

\section{Exact statement and proof}
Let $E_s$ be the full spectral ellipse and let $X$ be the declared closed candidate
rectangle. If $E_s\subseteq X$, then
\[
 E_s\cap X=E_s.
\]
This is the elementary set identity that intersection with a containing set leaves the
contained set unchanged.

\section{Source boundary}
This theorem \texttt{formalizes} the spectral containment premise in the corollary at
lines 737--773 of the Ragasa2026A v3 manuscript. Owner: \path{ksdft2effmass}; immutable
commit
\begin{center}\small\texttt{99a46cb59d0c6f27ada7d709f194996ef9619d9f}.\end{center}
Containment is an explicit hypothesis; this theorem does not verify numerical coordinate
extrema.

\section{Trusted base}
Lean and pinned Mathlib report \texttt{propext}, \texttt{Classical.choice}, and
\texttt{Quot.sound}. No admission or scientific axiom is used. The proof and prose were
AI-assisted and kernel-checked. No independent expert review is recorded. Status:
\texttt{proved}; working draft.
\end{document}
